\chapter{Encuesta aplicada a estudiantes de la asignatura FLP}
\label{annex:encuesta-estudiantes}

El siguiente instrumento fue aplicado a 27 estudiantes que cursaron la asignatura \textit{Fundamentos de Interpretación y Compilación de Lenguajes de Programación} en diferentes periodos académicos del programa de Ingeniería de Sistemas de la Universidad del Valle, sede Tuluá.

\section*{Instrucciones}

Por favor responda con sinceridad. Las respuestas son anónimas y tienen como único objetivo mejorar los recursos disponibles para los estudiantes de la asignatura.

\section*{Bloque I — Perfil académico}

\begin{enumerate}
    \item Semestre en que cursó la asignatura FLP:
    \begin{enumerate}
        \item Primero que la tomé.
        \item Segunda vez o más.
    \end{enumerate}

    \item Calificación final obtenida:
    \begin{enumerate}
        \item 3.0 – 3.4
        \item 3.5 – 3.9
        \item 4.0 – 4.4
        \item 4.5 – 5.0
        \item No la aprobé.
    \end{enumerate}
\end{enumerate}

\section*{Bloque II — Dificultades conceptuales}

Indique el nivel de dificultad que encontró en cada concepto durante el cursado de la asignatura, usando la siguiente escala:

\begin{center}
\textbf{1 = Sin dificultad \quad 2 = Poca dificultad \quad 3 = Dificultad moderada \quad 4 = Alta dificultad \quad 5 = Dificultad muy alta}
\end{center}

\begin{longtable}{|p{9cm}|p{0.5cm}|p{0.5cm}|p{0.5cm}|p{0.5cm}|p{0.5cm}|}
    \hline
    \textbf{Concepto} & \textbf{1} & \textbf{2} & \textbf{3} & \textbf{4} & \textbf{5} \\
    \hline
    Definición de gramáticas libres de contexto (BNF/EBNF) & & & & & \\ \hline
    Reglas de producción y derivaciones & & & & & \\ \hline
    Análisis sintáctico (\textit{parsing}) y sus algoritmos & & & & & \\ \hline
    Construcción de árboles de sintaxis abstracta (AST) & & & & & \\ \hline
    Interpretación del AST: recorrido y evaluación de nodos & & & & & \\ \hline
    Ambientes de ejecución y su representación & & & & & \\ \hline
    Alcance léxico (\textit{lexical scoping}) & & & & & \\ \hline
    Clausuras (\textit{closures}) y funciones de orden superior & & & & & \\ \hline
    Recursión en contextos de interpretación funcional & & & & & \\ \hline
    Definiciones dirigidas por sintaxis (SDD) y esquemas de traducción (SDT) & & & & & \\ \hline
    \caption{Instrumento — Dificultad conceptual percibida por los estudiantes.\\
    \scriptsize \textbf{Fuente:} Elaboración propia.}
    \label{tab:instrumento-dificultad}
\end{longtable}

\section*{Bloque III — Herramientas y recursos}

Indique su grado de acuerdo con cada afirmación (\textbf{Sí / No / Parcialmente}):

\begin{enumerate}
    \setcounter{enumi}{2}

    \item Racket fue una herramienta útil para comprender los conceptos de la asignatura.
    \begin{enumerate}
        \item Sí
        \item No
        \item Parcialmente
    \end{enumerate}

    \item La librería EOPL facilitó mi comprensión práctica del material del curso.
    \begin{enumerate}
        \item Sí
        \item No
        \item Parcialmente
    \end{enumerate}

    \item Encontré con facilidad recursos externos (tutoriales, videos, ejercicios) para complementar el estudio de la asignatura.
    \begin{enumerate}
        \item Sí
        \item No
        \item Parcialmente
    \end{enumerate}

    \item Me sentí seguro al aplicar los conceptos teóricos en los proyectos o evaluaciones de la asignatura.
    \begin{enumerate}
        \item Sí
        \item No
        \item Parcialmente
    \end{enumerate}
\end{enumerate}

\section*{Bloque IV — Recursos adicionales}

\begin{enumerate}
    \setcounter{enumi}{6}

    \item ¿Considera que contar con una herramienta web interactiva que le permitiera visualizar la construcción del AST y los cambios en el ambiente de ejecución habría mejorado su comprensión de la asignatura?
    \begin{enumerate}
        \item Sí, definitivamente.
        \item Probablemente sí.
        \item No lo sé.
        \item Probablemente no.
        \item No, no habría cambiado nada.
    \end{enumerate}

    \item ¿Qué funcionalidades consideraría más útiles en una herramienta de apoyo para la asignatura? \textit{(puede marcar varias)}:
    \begin{enumerate}
        \item Visualización gráfica del árbol de sintaxis abstracta (AST).
        \item Representación paso a paso de los cambios en el ambiente de ejecución.
        \item Editor de gramáticas con retroalimentación inmediata.
        \item Ejemplos predefinidos progresivos (de menor a mayor complejidad).
        \item Otra: \underline{\hspace{6cm}}
    \end{enumerate}

    \item Comentarios adicionales: \underline{\hspace{12cm}}
\end{enumerate}
